\documentclass[12pt]{article}
\usepackage[utf8]{inputenc}
\usepackage[T1]{fontenc}
\usepackage{hyperref}
\title{Wirangu-Nauo Language Resources for K-12 Curriculum Development}
\author{Compiled by Claude Sonnet 5.5 (claude-sonnet-5-5) for Andr\'as Kornai}
\date{}
\begin{document}
\maketitle

\section{Introduction}
The task gives the code wiw and the name Wirangu-Nauo. The family field says ``Bookkeeping'', which is a Glottolog label for unclassified entries and not a real family. My searches mostly surfaced resources for Wirangu of the far west coast of South Australia (Ceduna area), whose ISO code in the Wikipedia results is wgu. I could not confirm from the search results how wiw relates to that language, so everything below concerns Wirangu and should be checked before use. Wikipedia describes Wirangu as moribund, with about two semi-speakers reported in 2007: \url{https://en.wikipedia.org/wiki/Wirangu_language}.

\section{Primary Educational Resources}
Gladys Miller, one of the last two speakers, worked with linguists on a Wirangu picture dictionary, a \emph{Learner's Guide to Wirangu} and the children's book \emph{Wardugu Wirn} (a wombat hunt). The news item describing this is at \url{https://www.adelaide.edu.au/adelaidean/issues/12741/news12757.html}. I did not see download links for these materials.

\section{Linguistic and Reference Materials}
\emph{A Dictionary of the Wirangu Language of the Far West Coast of South Australia} (Hercus, Miller, Monaghan and Naessan; Tjutjunaku Worka Tjuta and the University of Adelaide, 2010, 118 pp.) has a catalogue record at \url{https://researchportalplus.anu.edu.au/en/publications/a-dictionary-of-the-wirangu-language-of-the-far-west-coast-of-sou/}.

The Wirangu Language Project dataset in the Aboriginal \&amp; Torres Strait Islander Data Archive has mixed access, with some material restricted: \url{https://www.atsida.edu.au/archive/datasets/wlp}.

\section{Religious and Community Resources}
The Mobile Language Team page may be relevant to community language work: \url{https://mobilelanguageteam.com.au/?p=227}. I did not verify its content.

\section{Digital and Audio Resources}
The roughly 200-word talking picture dictionary was distributed on CD, but I did not see an online version. The Ceduna Arts and Culture Centre has a language exhibition area where the Wirangu and Gugada dictionaries and language cards can be used.

\section{Curriculum Development Context}
A few published resources exist, namely the dictionary, a learner's guide and a children's book. They would have to be accessed through the community and the publishers. Any curriculum should be developed with the Wirangu community, and the identification of wiw with this material needs confirmation.
\end{document}